\subsubsection{Categorical constructions used by the evaluator}
\label{subsec:categorical-evaluator-constructions}

\begin{definition}[Twisted categorical input]
\label{def:twisted-categorical-input}
A twisted categorical input is a pair \((\cC,\eta)\), where \(\cC\) is a
premodular fusion category over \(\mathbb C\) and
\[
  \eta\in A_{\cC}
  =\operatorname{Aut}_{\otimes}(\operatorname{Id}_{\cC}).
\]
\end{definition}

\begin{definition}[Kirby colour and tensor-unit projector]
\label{def:kirby-colour}
\label{def:tensor-unit-projector}
Let \(\cC\) be a premodular fusion category.  Its Kirby colour is
\begin{equation}
  \Omega_{\cC}:=
  \sum_{i\in\operatorname{Irr}(\cC)}d_i\,i.
  \label{eq:kirby-colour}
\end{equation}
For an object \(W\), choose a basis
\(u_\mu\colon\mathbf1\to W\) and its composition-dual basis
\(\widehat u_\mu\colon W\to\mathbf1\), and set
\begin{equation}
  P_{\mathbf1,W}:=
  \sum_\mu u_\mu\widehat u_\mu
  \in\operatorname{End}_{\cC}(W).
  \label{eq:tensor-unit-projector}
\end{equation}
We also write
\(\operatorname{Cut}_{\cC}(W):=P_{\mathbf1,W}\); no global-dimension factor is
included in this notation.
\end{definition}

\begin{lemma}[Basis independence of the tensor-unit projector]
\label{lem:tensor-unit-projector-basis-independence}
Let \(\cC\) be a premodular fusion category and let \(W\in\cC\).
The endomorphism \(P_{\mathbf1,W}\) in
\eqref{eq:tensor-unit-projector} is independent of the chosen pair of dual
bases.
\end{lemma}

\begin{proof}
If \(v_\nu=\sum_\mu A_{\nu\mu}u_\mu\), then its dual basis is
\(\widehat v_\nu=\sum_\rho(A^{-1})_{\rho\nu}\widehat u_\rho\).  Hence
\[
  \sum_\nu v_\nu\widehat v_\nu
  =\sum_{\mu,\rho}
    \left(\sum_\nu A_{\nu\mu}(A^{-1})_{\rho\nu}\right)
    u_\mu\widehat u_\rho
  =\sum_\mu u_\mu\widehat u_\mu.
\]
\end{proof}

\begin{lemma}[Absorption by the tensor-unit projector]
\label{lem:unit-projector-absorbs-character}
Let \((\cC,\eta)\) be a twisted categorical input.
For every object $W\in\cC$ and every $n\in\mathbb Z$,
\[
  P_{\mathbf1,W}\eta_W^n=P_{\mathbf1,W}
  =\eta_W^nP_{\mathbf1,W}.
\]
\end{lemma}

\begin{proof}
Write $P_{\mathbf1,W}=\sum_\mu u_\mu\widehat u_\mu$ as in
\cref{def:tensor-unit-projector}.  Monoidality gives
$\eta_{\mathbf1}=\operatorname{id}_{\mathbf1}$.  Naturality for
$u_\mu\colon\mathbf1\to W$ and
$\widehat u_\mu\colon W\to\mathbf1$ gives
\[
  \eta_W^n u_\mu=u_\mu,
  \qquad
  \widehat u_\mu\eta_W^n=\widehat u_\mu.
\]
Summing proves both equalities.
\end{proof}

\begin{lemma}[Naturality of the tensor-unit projector]
\label{lem:unit-projector-naturality}
Let \(\cC\) be a premodular fusion category.
For every morphism $f\colon U\to V$ in $\cC$,
\[
  fP_{\mathbf1,U}=P_{\mathbf1,V}f.
\]
\end{lemma}

\begin{proof}
Finite semisimplicity decomposes every object into simple isotypic summands.
By Schur's lemma, a morphism has no matrix entries between the tensor-unit
isotypic summand and any nonisomorphic simple isotypic summand.  It therefore
preserves the tensor-unit summand.
\end{proof}

\begin{corollary}[Finite transport of the tensor-unit projector]
\label{cor:finite-unit-projector-transport}
Let \(\cC\) be a premodular fusion category.
A tensor-unit projector may be moved through any finite composite of
crossings, cups, caps, twists, and morphism coupons, with the source and
target projectors changed according to the boundary objects.
\end{corollary}

\begin{proof}
Apply \cref{lem:unit-projector-naturality} successively to the elementary
morphisms in the ribbon expression.  Graphically, this is the tunnelling
property of the completeness idempotent in
\cite[Lemmas~2.27--2.28]{Tham2021}.
\end{proof}

\begin{lemma}[Character transport along a ribbon]
\label{lem:finite-coupon-transport}
Let \((\cC,\eta)\) be a twisted categorical input and let
$n\in\mathbb Z$.
\begin{enumerate}[label=(\roman*),ref=\roman*]
  \item For every morphism $f\colon U\to V$ in \(\cC\),
  \[
    (\eta^n)_Vf=f(\eta^n)_U.
  \]

  \item The same naturality identity applies to braidings and twists.

  \item For every object \(X\in\cC\), at a dual bend the transported
  component is its mate:
  \[
    \operatorname{ev}_X
      (\operatorname{id}_{X^*}\otimes(\eta^n)_X)
    =
    \operatorname{ev}_X
      ((\eta^{-n})_{X^*}\otimes\operatorname{id}_X),
  \]
  with the analogous coevaluation identity.
\end{enumerate}
Consequently a character decoration may be transported along any compact
ribbon through the finitely many crossings, extrema, twists, and coupons of
a generic diagram.
\end{lemma}

\begin{proof}
The first statement is naturality of the monoidal natural automorphism
$\eta^n$; braiding and twist are particular morphisms.  Naturality for
$\operatorname{ev}_X$ gives
\[
 \operatorname{ev}_X
 ((\eta^n)_{X^*}\otimes(\eta^n)_X)=\operatorname{ev}_X.
\]
Precompose with $(\eta^{-n})_{X^*}\otimes\operatorname{id}_X$ to obtain the
displayed bend identity; the coevaluation formula is dual.  Applying these
identities successively along the finite route proves the final assertion.
\end{proof}

\begin{remark}[Scope of coupon naturality]
\label{rem:coupon-naturality-scope}
The first identity transports the automorphism of the whole boundary object
through a multivalent coupon; it does not route one factor through an
arbitrary leg.  At a dual bend, expressing the same oriented route on the
dual strand replaces the local exponent \(n\) by \(-n\).
\end{remark}

\begin{lemma}[Two categorical cancellation sums]
\label{lem:local-cancellation-sums}
Let \(\cC\) be a premodular fusion category.
For the simple-label set of $\cC$,
\begin{equation}
  \sum_i d_i\,\operatorname{tr}_i(P_{\mathbf1,i})=1,
  \qquad
  \sum_i d_i\,\operatorname{tr}_i(\operatorname{id}_i)
  =\operatorname{Dim}(\cC).
  \label{eq:local-cancellation-sums}
\end{equation}
\end{lemma}

\begin{proof}
If $i$ is simple, then $P_{\mathbf1,i}$ is the identity for
$i\cong\mathbf1$ and zero otherwise.  Since $d_{\mathbf1}=1$, the first sum
is $1$.  Spherical trace gives
$\operatorname{tr}_i(\operatorname{id}_i)=d_i$, so the second sum is
$\sum_i d_i^2=\operatorname{Dim}(\cC)$.
\end{proof}
